% !TeX program = xelatex
% 中文 + 专业数学 + 参考文献 的最小可编译模板（2026-09-15 在本机实测通过）
% 编译：xelatex -> biber -> xelatex -> xelatex   （或 python build.py template.tex）
\documentclass[11pt,a4paper]{ctexart}

% ---------- 字体：中文由 ctexart 处理；西文与数学都用 TeX Gyre Termes 一套 ----------
% 前置条件：miktex packages install tex-gyre tex-gyre-math
% 若不想装包，把下面两行换成 Cambria / Cambria Math（Windows 自带真 MATH 字体）
\usepackage{unicode-math}
\setmainfont{TeX Gyre Termes}
\setmathfont{TeX Gyre Termes Math}

\usepackage{amsmath}
\usepackage{booktabs}
\usepackage{amsthm}
\usepackage[backend=biber,style=numeric]{biblatex}
\addbibresource{refs.bib}

% ---------- 中文本地化：默认是英文，漏一个就露出来 ----------
\renewcommand{\contentsname}{目录}
\renewcommand{\figurename}{图}
\renewcommand{\tablename}{表}
\renewcommand{\proofname}{证明}
% 参考文献标题：实测 \renewcommand{\refname} 与 \DefineBibliographyStrings
% **都改不动 biblatex 的标题**，唯一可靠的是显式定义 heading：
\defbibheading{bibliography}[\refname]{\section*{参考文献}}

% 题注分隔符：中文体例不用半角冒号，改成空格
\usepackage{caption}
\captionsetup{labelsep=quad}

\theoremstyle{definition}
\newtheorem{theorem}{定理}

% 若要求"行尾不悬算符"，打开下面两行（代价：长公式可能溢出）
% \binoppenalty=10000
% \relpenalty=10000

\pagestyle{plain}

\begin{document}

\section*{最小模板}

行内根号：$\sqrt{2}$、$\sqrt[3]{2}$、$\sqrt{1+\sqrt{2}}$、$\sqrt{\dfrac{p}{q}}$。

GB/T 3102.11 正斜体：变量 $x$ 斜体；常数 $\symup{\pi}$、$\symup{e}$、$\symup{i}^{2}=-1$ 正体；
函数名 $\sin$、$\exp$、$\ln$ 正体；算子 $\mathrm{d}f/\mathrm{d}x$ 正体。

$$\sum_{n=1}^{\infty}\frac{1}{n^{2}}=\frac{\symup{\pi}^{2}}{6}
\qquad
\lim_{\delta\to0}(1+\delta)^{1/\delta}=\symup{e}$$

\begin{theorem}[示例]
  不存在正整数 $p$、$q$ 使 $(p/q)^{2}=2$。
\end{theorem}
\begin{proof}
  设 $p^{2}=2q^{2}$ 且 $\gcd(p,q)=1$，推出 $p$、$q$ 同为偶数，与互质矛盾。
\end{proof}

\begin{table}[htbp]
  \centering
  \begin{tabular}{lrr}
    \toprule
    角差 $\delta$ & 深度 $h$ & 折算 \\
    \midrule
    $2'$  & $1.72$ 公里  & --- \\
    $2''$ & $103.1$ 公里 & --- \\
    \bottomrule
  \end{tabular}
  \caption{数值区间用 en dash：约公元前 1800--前 1600 年}
\end{table}

引用一处文献\cite{example2026}。

\printbibliography

\end{document}
